% ============================================================
% TIKZ DRAWING LIBRARY — Đồ thị, Bảng biến thiên, Hình học
% ============================================================
% File: tikz/lib/tikz-math.tex
% Mô tả: Tập hợp các command & environment TikZ chuyên biệt
%         cho bài toán THPT QG (Toán 12)
% ============================================================

% ============================================================
% PHẦN 1: BẢNG BIẾN THIÊN (BBT)
% ============================================================
%
% Cú pháp:
%   \bangbienthien{
%     xvals = {x_1, x_2, ..., x_n},        % các giá trị x đặc biệt
%     fprimevals = {+, 0, -, 0, +, ...},    % dấu f' tương ứng các khoảng
%     fvals = {y_1, y_2, ..., y_n},         % giá trị f(x) tại các x
%     arrows = {up, down, up, ...},         % chiều mũi tên (up/down)
%   }
%
% Cách hoạt động:
%   - Dùng tikz table thủ công, linh hoạt hoàn toàn
%   - Hỗ trợ $-\infty$ và $+\infty$ ở đầu/cuối
% ============================================================

% === Style cho BBT ===
\tikzset{
  bbt/.style={
    draw=ForumDarkBlue,
    line width=0.6pt,
    font=\small,
  },
  bbt header/.style={
    fill=ForumDarkBlue!10,
    font=\small\bfseries\color{ForumDarkBlue},
  },
  bbt arrow/.style={
    ->,
    >=stealth,
    line width=0.7pt,
    color=ForumMidBlue,
  },
  bbt val/.style={
    font=\small,
    text=ForumDarkGray,
  },
  bbt sign/.style={
    font=\small,
    text=ForumDarkBlue,
  },
}

% --- BBT đơn giản: 2 điểm cực trị (hàm bậc 3) ---
% \bbtBacBa{x1}{x2}{f(x1)}{f(x2)}{lim-inf}{lim+inf}
\newcommand{\bbtBacBa}[6]{%
\begin{center}
\begin{tikzpicture}[scale=1, every node/.style={inner sep=3pt}]
  % Khung ngoài
  \draw[bbt] (0,0) rectangle (10,-4.5);
  
  % Hàng header: x
  \draw[bbt] (0,-1) -- (10,-1);
  \draw[bbt] (1.5,0) -- (1.5,-4.5);
  
  % Header labels
  \node[bbt header] at (0.75, -0.5) {$x$};
  \node[bbt header] at (0.75, -1.75) {$f'(x)$};
  \draw[bbt] (0,-2.5) -- (10,-2.5);
  \node[bbt header] at (0.75, -3.5) {$f(x)$};
  
  % x values
  \node[bbt val] at (2.2, -0.5) {$-\infty$};
  \node[bbt val] at (4.5, -0.5) {$#1$};
  \draw[bbt, dashed, opacity=0.4] (4.5,-1) -- (4.5,-4.5);
  \node[bbt val] at (7, -0.5) {$#2$};
  \draw[bbt, dashed, opacity=0.4] (7,-1) -- (7,-4.5);
  \node[bbt val] at (9.2, -0.5) {$+\infty$};
  
  % f'(x) signs
  \node[bbt sign] at (3.3, -1.75) {$+$};
  \node[bbt sign] at (4.5, -1.75) {$0$};
  \node[bbt sign] at (5.75, -1.75) {$-$};
  \node[bbt sign] at (7, -1.75) {$0$};
  \node[bbt sign] at (8.1, -1.75) {$+$};
  
  % f(x) values and arrows
  \node[bbt val] at (2.2, -3.8) {$#5$};
  \node[bbt val] at (4.5, -2.8) {$#3$};
  \node[bbt val] at (7, -4.2) {$#4$};
  \node[bbt val] at (9.2, -2.8) {$#6$};
  
  % Arrows
  \draw[bbt arrow] (2.5,-3.6) -- (4.2,-3.0);
  \draw[bbt arrow] (4.8,-3.0) -- (6.7,-4.0);
  \draw[bbt arrow] (7.3,-4.0) -- (8.9,-3.0);
\end{tikzpicture}
\end{center}
}

% --- BBT hàm bậc 4: dạng y = ax^4 + bx^2 + c ---
% \bbtBacBon{x1}{x2}{x3}{f(x1)}{f(x2)}{f(x3)}{lim-inf}{lim+inf}
\newcommand{\bbtBacBon}[8]{%
\begin{center}
\begin{tikzpicture}[scale=1, every node/.style={inner sep=3pt}]
  \draw[bbt] (0,0) rectangle (12,-4.5);
  \draw[bbt] (0,-1) -- (12,-1);
  \draw[bbt] (1.5,0) -- (1.5,-4.5);
  \draw[bbt] (0,-2.5) -- (12,-2.5);
  
  % Headers
  \node[bbt header] at (0.75, -0.5) {$x$};
  \node[bbt header] at (0.75, -1.75) {$f'(x)$};
  \node[bbt header] at (0.75, -3.5) {$f(x)$};
  
  % x values
  \node[bbt val] at (2.2, -0.5) {$-\infty$};
  \node[bbt val] at (4, -0.5) {$#1$};
  \draw[bbt, dashed, opacity=0.4] (4,-1) -- (4,-4.5);
  \node[bbt val] at (6, -0.5) {$#2$};
  \draw[bbt, dashed, opacity=0.4] (6,-1) -- (6,-4.5);
  \node[bbt val] at (8, -0.5) {$#3$};
  \draw[bbt, dashed, opacity=0.4] (8,-1) -- (8,-4.5);
  \node[bbt val] at (10.5, -0.5) {$+\infty$};
  
  % f'(x)
  \node[bbt sign] at (3.1, -1.75) {$-$};
  \node[bbt sign] at (4, -1.75) {$0$};
  \node[bbt sign] at (5, -1.75) {$+$};
  \node[bbt sign] at (6, -1.75) {$0$};
  \node[bbt sign] at (7, -1.75) {$-$};
  \node[bbt sign] at (8, -1.75) {$0$};
  \node[bbt sign] at (9.2, -1.75) {$+$};
  
  % f(x) values and arrows
  \node[bbt val] at (2.2, -2.8) {$#7$};
  \node[bbt val] at (4, -4.2) {$#4$};
  \node[bbt val] at (6, -2.8) {$#5$};
  \node[bbt val] at (8, -4.2) {$#6$};
  \node[bbt val] at (10.5, -2.8) {$#8$};
  
  \draw[bbt arrow] (2.5,-3.0) -- (3.7,-4.0);
  \draw[bbt arrow] (4.3,-4.0) -- (5.7,-3.0);
  \draw[bbt arrow] (6.3,-3.0) -- (7.7,-4.0);
  \draw[bbt arrow] (8.3,-4.0) -- (10.2,-3.0);
\end{tikzpicture}
\end{center}
}

% --- BBT tổng quát (custom) ---
% Môi trường vẽ BBT tự do dùng khi các macro trên không đủ
% Sử dụng: \begin{bangBT}{width}{height} ... \end{bangBT}
\newenvironment{bangBT}[2]{%
  \begin{center}
  \begin{tikzpicture}[scale=1, every node/.style={inner sep=3pt}]
  \draw[bbt] (0,0) rectangle (#1,#2);
}{%
  \end{tikzpicture}
  \end{center}
}

% --- BBT tối ưu: 1 điểm cực trị trên miền ---
% \bbtOptim[var]{a}{x0}{b}{f(a)}{f(x0)}{f(b)}{fname}{type}
%   var   = tên biến (mặc định x)
%   a, b  = biên miền (có thể dùng -\infty, +\infty)
%   x0    = điểm cực trị
%   type  = max hoặc min
\newcommand{\bbtOptim}[9][x]{%
\begin{center}
\begin{tikzpicture}[scale=1, every node/.style={inner sep=3pt}]
  % Frame
  \draw[bbt] (0,0) rectangle (9,-4.2);
  \draw[bbt] (0,-1) -- (9,-1);
  \draw[bbt] (0,-2.2) -- (9,-2.2);
  \draw[bbt] (1.4,0) -- (1.4,-4.2);
  % Headers
  \node[bbt header] at (0.7,-0.5) {$#1$};
  \node[bbt header] at (0.7,-1.6) {$#8'$};
  \node[bbt header] at (0.7,-3.2) {$#8$};
  % x values
  \node[bbt val] at (2.3,-0.5) {$#2$};
  \node[bbt val] at (5,-0.5) {$#3$};
  \draw[bbt,dashed,opacity=0.4] (5,-1) -- (5,-4.2);
  \node[bbt val] at (7.8,-0.5) {$#4$};
  % Choose max or min layout
  \def\bbttype{#9}%
  \def\bbtmax{max}%
  \ifx\bbttype\bbtmax
    % f' : + 0 -
    \node[bbt sign] at (3.5,-1.6) {$+$};
    \node[bbt sign] at (5,-1.6) {$0$};
    \node[bbt sign] at (6.5,-1.6) {$-$};
    % f values: low → high → low
    \node[bbt val] at (2.3,-3.8) {$#5$};
    \node[bbt val] at (5,-2.6) {$#6$};
    \node[bbt val] at (7.8,-3.8) {$#7$};
    % arrows up then down
    \draw[bbt arrow] (2.6,-3.6) -- (4.7,-2.8);
    \draw[bbt arrow] (5.3,-2.8) -- (7.5,-3.6);
  \else
    % f' : - 0 +
    \node[bbt sign] at (3.5,-1.6) {$-$};
    \node[bbt sign] at (5,-1.6) {$0$};
    \node[bbt sign] at (6.5,-1.6) {$+$};
    % f values: high → low → high
    \node[bbt val] at (2.3,-2.6) {$#5$};
    \node[bbt val] at (5,-3.8) {$#6$};
    \node[bbt val] at (7.8,-2.6) {$#7$};
    % arrows down then up
    \draw[bbt arrow] (2.6,-2.8) -- (4.7,-3.6);
    \draw[bbt arrow] (5.3,-3.6) -- (7.5,-2.8);
  \fi
\end{tikzpicture}
\end{center}
}


% ============================================================
% PHẦN 2: ĐỒ THỊ HÀM SỐ
% ============================================================

% === Pgfplots global settings ===
\pgfplotsset{
  compat=1.18,
  % Style chuẩn cho đồ thị hàm số
  forum graph/.style={
    axis lines=middle,
    xlabel={$x$},
    ylabel={$y$},
    xlabel style={at={(ticklabel* cs:1)}, anchor=west, font=\small},
    ylabel style={at={(ticklabel* cs:1)}, anchor=south, font=\small},
    tick label style={font=\footnotesize},
    grid=major,
    grid style={TikzGrid, line width=0.3pt},
    axis line style={TikzAxis, line width=0.8pt, -stealth},
    width=9cm,
    height=7cm,
    samples=200,
    every axis plot/.append style={line width=1.2pt},
    legend style={
      at={(0.98,0.98)},
      anchor=north east,
      font=\footnotesize,
      draw=ForumGray,
      fill=white,
      fill opacity=0.9,
    },
  },
  % Style nhỏ gọn (dùng trong lời giải)
  forum graph small/.style={
    forum graph,
    width=7cm,
    height=5.5cm,
  },
  % Style cho đồ thị tham số
  forum param/.style={
    forum graph,
    axis equal,
  },
}

% --- Vẽ đường tiệm cận ---
% \tiemcan{x/y}{giá trị}{xmin}{xmax} hoặc {ymin}{ymax}
\newcommand{\tiemcanDung}[3]{%
  \draw[dashed, TikzAsymptote, line width=0.8pt] 
    (axis cs:#1,#2) -- (axis cs:#1,#3);
}
\newcommand{\tiemcanNgang}[3]{%
  \draw[dashed, TikzAsymptote, line width=0.8pt] 
    (axis cs:#2,#1) -- (axis cs:#3,#1);
}

% --- Đánh dấu điểm đặc biệt ---
% \diemDB{x}{y}{label}{vị trí label}
\newcommand{\diemDB}[4]{%
  \node[circle, fill=ForumMidBlue, inner sep=1.8pt] at (axis cs:#1,#2) {};
  \node[font=\footnotesize\color{ForumDarkBlue}, #4] at (axis cs:#1,#2) {#3};
}

% --- Đánh dấu điểm rỗng (hở) ---
\newcommand{\diemHo}[4]{%
  \node[circle, draw=ForumMidBlue, fill=white, inner sep=1.5pt, line width=0.8pt] 
    at (axis cs:#1,#2) {};
  \node[font=\footnotesize\color{ForumDarkBlue}, #4] at (axis cs:#1,#2) {#3};
}

% --- Tô vùng (dùng cho tích phân) ---
% Sử dụng trực tiếp \addplot[fill=TikzFill, opacity=0.3] ...

% ============================================================
% PHẦN 3: HÌNH HỌC KHÔNG GIAN
% ============================================================

\tikzset{
  % Cạnh thấy
  solid edge/.style={
    line width=0.8pt,
    color=ForumDarkBlue,
  },
  % Cạnh khuất 
  hidden edge/.style={
    line width=0.5pt,
    dashed,
    color=ForumGray,
  },
  % Đỉnh
  vertex/.style={
    circle,
    fill=ForumDarkBlue,
    inner sep=1.5pt,
  },
  % Label đỉnh
  vertex label/.style={
    font=\small\bfseries\color{ForumDarkBlue},
  },
  % Góc vuông
  right angle mark/.style={
    line width=0.5pt,
    color=ForumMidBlue,
  },
  % Mặt phẳng tô
  plane fill/.style={
    fill=TikzFill,
    fill opacity=0.2,
    draw=ForumMidBlue,
    draw opacity=0.5,
  },
  % Chiều cao / khoảng cách
  height line/.style={
    line width=0.6pt,
    color=ForumRed,
    dashed,
  },
}

% --- Vẽ ký hiệu góc vuông ---
% \gocVuong{điểm đỉnh}{điểm 1}{điểm 2}{kích thước}
\newcommand{\gocVuong}[4]{%
  \coordinate (tmp1) at ($(#1)!#4!(#2)$);
  \coordinate (tmp2) at ($(#1)!#4!(#3)$);
  \coordinate (tmp3) at ($(tmp1)+(tmp2)-(#1)$);
  \draw[right angle mark] (tmp1) -- (tmp3) -- (tmp2);
}

% --- Vẽ đỉnh + label ---
\newcommand{\dinh}[3]{%
  \node[vertex] (#1) at #2 {};
  \node[vertex label, #3] at #2 {$#1$};
}


% ============================================================
% PHẦN 4: SỐ PHỨC — MẶT PHẲNG PHỨC
% ============================================================

\tikzset{
  complex plane/.style={
    axis lines=middle,
    xlabel={$\text{Re}$},
    ylabel={$\text{Im}$},
  },
  complex point/.style={
    circle,
    fill=ForumMidBlue,
    inner sep=1.8pt,
  },
  complex circle/.style={
    draw=ForumMidBlue,
    line width=0.8pt,
  },
  complex region/.style={
    fill=TikzFill,
    fill opacity=0.3,
    draw=ForumMidBlue,
    line width=0.6pt,
  },
}

% ============================================================
% PHẦN 5: XÁC SUẤT — CÂY XÁC SUẤT
% ============================================================

\tikzset{
  prob node/.style={
    circle,
    draw=ForumDarkBlue,
    fill=ForumLightBlue!50,
    minimum size=8mm,
    font=\small,
    line width=0.6pt,
  },
  prob edge/.style={
    ->,
    >=stealth,
    line width=0.6pt,
    color=ForumDarkBlue,
  },
  prob label/.style={
    font=\footnotesize,
    color=ForumOrange,
    midway,
  },
  prob result/.style={
    font=\small\bfseries,
    color=ForumDarkBlue,
  },
}
